\section{Principle}
\paragraph{Purpose}
The knowledge of an LLM is frozen in its weights at the end of training: it is
\tB{outdated}, it does not contain private documents, and the model can state false facts
with confidence (\emph{hallucination}). \textbf{Retrieval-Augmented Generation} (RAG, Lewis et
al., 2020) \uB{retrieves the relevant passages at query time and puts them in the prompt}, so
that the model answers from them. The knowledge is changed by updating the index, with no
retraining, and the sources can be cited.

\paragraph{Pipeline}
\begin{enumerate}
    \item \textbf{indexing} (offline): cut the documents into \emph{chunks}, compute an
        embedding vector for each chunk, store them in a vector index;
    \item \textbf{retrieval} (online): embed the query, find the $k$ nearest chunks;
    \item \textbf{generation}: build a prompt with the instruction, the retrieved chunks and
        the question, and let the model answer.
\end{enumerate}

\section{Embeddings for Retrieval}
\paragraph{Similarity}
A text embedding model (a Transformer encoder, with pooling of the token vectors into one
vector) maps a query $\bm{q}$ and a chunk $\bm{d}$ to $\mathbb{R}^{m}$, and their relevance
is the \textbf{cosine similarity}
\begin{center}
    $\cos(\bm{q},\bm{d})=\dfrac{\bm{q}^{T}\bm{d}}{\norm{\bm{q}}\,\norm{\bm{d}}}$
\end{center}
which equals the dot product when the vectors are normalized.

\paragraph{Contrastive training}
The model is trained so that a query is close to its relevant chunk $d^{+}$ and far from the
others. With a batch of pairs, the other chunks of the batch are used as negatives
(\emph{InfoNCE}) and $\tau$ is a temperature:
\begin{center}
    \frB{$\mathcal{L}=-\log\dfrac{\exp\left(\cos(\bm{q},\bm{d}^{+})/\tau\right)}{\su{j}{}
    \exp\left(\cos(\bm{q},\bm{d}_{j})/\tau\right)}$}
\end{center}
It is the softmax cross-entropy where the classes are the chunks of the batch.

\section{Vector Search}
Exact search compares the query with the $N$ vectors, in $O(Nm)$: too slow for millions of
chunks. \textbf{Approximate nearest-neighbor} (ANN) indexes trade a little recall for a large
speed-up:
\begin{itemize}
    \item \textbf{IVF}: partition the vectors with $k$-means, and search only the few
        clusters closest to the query;
    \item \textbf{HNSW}: a multi-layer proximity graph, the search walks greedily from
        coarse to fine layers;
    \item \textbf{product quantization}: compress the vectors to save memory.
\end{itemize}

\section{Chunking}
The size of the chunks is a trade-off:
\begin{itemize}
    \item \tB{too small}: a chunk loses the context needed to be understood, for example the
        subject of a sentence;
    \item \tB{too large}: the embedding averages several topics and is less precise, and the
        chunks use up the context window.
\end{itemize}
Typical sizes are a few hundred tokens, with an \emph{overlap} between consecutive chunks, or
a cut at the structure of the document (sections, paragraphs).

\section{Hybrid Search and Reranking}
\paragraph{BM25 (lexical)} Dense embeddings capture meaning but can miss exact terms (names,
identifiers, codes). BM25 scores the words shared by the query and the document $d$:
\begin{center}
    $\mathrm{BM25}(q,d)=\su{t\in q}{}IDF(t)\cdot\dfrac{f(t,d)\,(k_{1}+1)}{f(t,d)+k_{1}
    \left(1-b+b\frac{|d|}{avgdl}\right)}$
\end{center}
$f(t,d)$ is the frequency of the term $t$ in $d$, $IDF$ gives more weight to rare terms,
and the denominator saturates the frequency and normalizes by the length ($k_{1}\approx1.2$ to
$2$, $b=0.75$).

\paragraph{Hybrid} Run both searches and merge the rankings, for example with
\emph{reciprocal rank fusion}: $score(d)=\su{i}{}\dfrac{1}{c+rank_{i}(d)}$ ($c\approx60$),
which needs no calibration between the different scores.

\paragraph{Reranking} A \textbf{cross-encoder} reads the query and a chunk \uB{together} in
a Transformer and outputs a relevance score. It is more accurate than comparing two
independent embeddings but too slow for the whole collection, so it is applied to the top
$100$ of the first stage to keep the best few.

\section{Evaluation}
\begin{itemize}
    \item \textbf{retrieval}: $recall@k$ (share of the relevant chunks found in the top $k$),
        $precision@k$, MRR and nDCG for the quality of the ranking;
    \item \textbf{generation}: \emph{faithfulness} (is the answer supported by the retrieved
        chunks?) and \emph{relevance} (does it answer the question?).
\end{itemize}
The two stages must be evaluated separately, to know whether a wrong answer comes from a
retrieval miss or from the generation.

\section{Failure Modes}
\begin{itemize}
    \item the relevant chunk is \tB{not retrieved} (bad chunking, vocabulary mismatch);
    \item irrelevant chunks \uB{distract} the model, which may follow them;
    \item \emph{lost in the middle}: models use information at the beginning and at the end
        of a long context better than in the middle, so the order of the chunks matters;
    \item \textbf{prompt injection}: a retrieved document can contain instructions, which the
        model may treat as if they came from the user.
\end{itemize}

\section{RAG, Fine-tuning and Long Context}
\begin{itemize}
    \item \textbf{RAG}: for facts that change or are private, with citations. The model does
        not learn a new style or skill;
    \item \textbf{fine-tuning}: for behavior, format and style, not for storing many facts;
    \item \textbf{long context}: put all the documents in the prompt when they fit, without
        an index, at a cost that grows with the length (attention is $O(n^{2})$) and with the
        risk of lost-in-the-middle.
\end{itemize}
